\subsection{\texorpdfstring{Example 1: the complex \(\mathbf{S}(\mathbf{L})\otimes_{\mathbf{A}}\boldsymbol{\Lambda}(\mathbf{L})\)}{Example 1: the complex \textbackslash mathbf\{S\}(\textbackslash mathbf\{L\})\textbackslash otimes\_\{\textbackslash mathbf\{A\}\}\textbackslash boldsymbol\{\textbackslash Lambda\}(\textbackslash mathbf\{L\})}}

Let A be a ring, L an A-module, \(\mathbf{S}(\mathrm{L})\) its symmetric algebra, \(\mathbf{S}(\mathrm{L}) \otimes_{\mathrm{A}} \mathrm{L}\) be the \(\mathbf{S}(\mathrm{L})\) -module obtained by extension of scalars, \(u: \mathbf{S}(\mathrm{L}) \otimes_{\mathrm{A}} \mathrm{L} \to \mathbf{S}(\mathrm{L})\) the linear form such that \(u(s \otimes x) = sx\) for \(s \in \mathbf{S}(\mathrm{L})\) , \(x \in \mathrm{L}\) . By the canonical isomorphism of \(\mathbf{S}(\mathrm{L})\) -modules (III, p.~83, prop. 8)

\[
\boldsymbol {\Lambda} _ {\mathbf {S} (L)} (\mathbf {S} (L) \otimes_ {A} L) \rightarrow \mathbf {S} (L) \otimes_ {A} \boldsymbol {\Lambda} (L),
\]

the differential of the complex \(\mathbf{K}^{\mathbf{S}(\mathrm{L})}(u)\) is transported to the map

\[
d: \mathbf {S} (\mathrm{L}) \otimes_ {\mathrm{A}} \boldsymbol {\Lambda} (\mathrm{L}) \rightarrow \mathbf {S} (\mathrm{L}) \otimes_ {\mathrm{A}} \boldsymbol {\Lambda} (\mathrm{L})
\]

such that, for \(x_{1},\ldots,x_{p},y_{1},\ldots,y_{q}\) in L, we have

\[
\begin{array}{l} d \left(\left(x _ {1} \dots x _ {p}\right) \otimes \left(y _ {1} \wedge \dots \wedge y _ {q}\right)\right) \\ = \sum_ {i = 1} ^ {q} (- 1) ^ {i + 1} y _ {i} x _ {1} \dots x _ {p} \otimes \left(y _ {1} \wedge \dots \wedge y _ {i - 1} \wedge y _ {i + 1} \wedge \dots \wedge y _ {q}\right). \end{array}
\]

Note that \(d\) maps \(\mathbf{S}^{p}(\mathrm{L})\otimes\symbf{\Lambda}^{q}(\mathrm{L})\) to \(\mathbf{S}^{p+1}(\mathrm{L})\otimes\symbf{\Lambda}^{q-1}(\mathrm{L})\) , so that the complex of A-modules \(\mathbf{S}(\mathrm{L})\otimes\symbf{\Lambda}(\mathrm{L})\) decomposes as the direct sum of the complexes described by the following diagrams:

\[
\left(\mathcal {E} _ {n}\right): 0 \rightarrow \mathbf {S} ^ {0} L \otimes_ {\mathrm{A}} \boldsymbol {\Lambda} ^ {n} L \rightarrow \mathbf {S} ^ {1} L \otimes_ {\mathrm{A}} \boldsymbol {\Lambda} ^ {n - 1} L \rightarrow \dots \rightarrow \mathbf {S} ^ {n} L \otimes_ {\mathrm{A}} \boldsymbol {\Lambda} ^ {0} L \rightarrow 0, n \in \mathbf {N}.
\]

If the A-module L is the direct sum of a finite family \((\mathbf{L}_{i})_{i\in I}\) where I is totally ordered, the canonical bijection

\[
\bigotimes_ {i \in I} ^ {\varepsilon} \left(\mathbf {S} (L _ {i}) \otimes_ {A} \boldsymbol {\Lambda} (L _ {i})\right) \to \mathbf {S} (L) \otimes_ {A} \boldsymbol {\Lambda} (L)
\]

is an isomorphism of complexes of A-modules (this follows from prop. 2 of X, p.~148 or from formula (9) above).

PROPOSITION 3. --- If the A-module L is flat, then the sequences \((\mathcal{E}_n)\) above are exact for \(n > 0\) .

\begin{enumerate}
\def\labelenumi{\alph{enumi})}
\tightlist
\item
  Let us first note that, if \(p_{L}\) is the composite homomorphism
\end{enumerate}

\[
\mathbf {S} (\mathrm{L}) \otimes \boldsymbol {\Lambda} (\mathrm{L}) \xrightarrow {\alpha} \mathbf {S} ^ {0} (\mathrm{L}) \otimes \boldsymbol {\Lambda} ^ {0} (\mathrm{L}) \xrightarrow {\beta} \mathrm{A},
\]

where \(\alpha\) is the tensor product of the canonical projections and \(\beta\) the canonical isomorphism; it remains to prove that \(\mathbf{H}(p_{\mathrm{L}})\) is bijective.

\begin{enumerate}
\def\labelenumi{\alph{enumi})}
\setcounter{enumi}{1}
\item
  If L = 0 or if L = A, the proposition is evident.
\item
  Suppose L is free of finite rank; write it as a direct sum \(L_{1} \oplus \ldots \oplus L_{n}\) of free A-modules of rank 1. By the remark preceding the proposition, the complex \(\mathbf{S}(\mathrm{L}) \otimes \symbf{\Lambda}(\mathrm{L})\) is isomorphic to the tensor product of the n free complexes \(\mathbf{S}(\mathrm{L}_{i}) \otimes \symbf{\Lambda}(\mathrm{L}_{i})\) whose homology is free by b). By X, p.~79, cor. 4, the canonical homomorphism
\end{enumerate}

\[
\gamma : \bigoplus_ {i = 1} ^ {n} \mathrm{H} (\mathbf {S} (\mathrm{L} _ {i}) \otimes \boldsymbol {\Lambda} (\mathrm{L} _ {i})) \rightarrow \mathrm{H} (\mathbf {S} (\mathrm{L}) \otimes \boldsymbol {\Lambda} (\mathrm{L}))
\]

is bijective. By b) \(\mathrm{H}(p_{\mathrm{L}_i})\) is bijective for all \(i\) . Since \(\bigotimes_{i=1}^{n} \mathrm{H}(p_{\mathrm{L}_i}) = \mathrm{H}(p_{\mathrm{L}}) \circ \gamma\) ,

\[
\mathrm{H} (p _ {\mathrm{L}}) \text {   is   bijective }  .
\]

\begin{enumerate}
\def\labelenumi{\alph{enumi})}
\setcounter{enumi}{3}
\tightlist
\item
  In the general case, L is the inductive limit of a filtered inductive system \((\mathbf{L}_{i})_{i\in I}\) of free modules of finite rank (X, p.~14, th. 1). Since the canonical bijective homomorphism
\end{enumerate}

\[
\varinjlim \mathbf {S} (\mathrm{L} _ {i}) \otimes \boldsymbol {\Lambda} (\mathrm{L} _ {i}) \rightarrow \mathbf {S} (\mathrm{L}) \otimes \boldsymbol {\Lambda} (\mathrm{L})
\]

is an isomorphism of complexes, the proposition follows from X, p.~28, prop. 1.

Remarks. --- 1) We shall see below (X, p.~158, example) another proof of part c) above.

\begin{enumerate}
\def\labelenumi{\arabic{enumi})}
\setcounter{enumi}{1}
\tightlist
\item
  If A is a Q-algebra, the conclusion of Prop. 3 remains true without any hypothesis on L (cf. X, p. 206, Exercise 1).
\end{enumerate}

Let G be a group and \(\rho: G \to \mathbf{GL}(L)\) a linear representation of G in a flat A-module L. Then the \((\mathcal{E}_n)\) are exact sequences of linear representations. Suppose L is finitely generated projective, and denote by \(R_A(G)\) the ring of representations of G in finitely generated projective A-modules. It follows from Prop. 3 that we have in \(R_{A}(G)\) the relations

\[
\sum_ {i = 0} ^ {n} (- 1) ^ {i} \left[ \mathbf {S} ^ {i} (\mathrm{L}) \right] \left[ \boldsymbol {\Lambda} ^ {n - i} (\mathrm{L}) \right] = 0, \quad n > 0.\tag{10}
\]

If we consider the formal series

\[
s (\mathrm{T}) = \sum_ {i = 0} ^ {\infty} \left[ \mathbf {S} ^ {i} (\mathrm{L}) \right] \mathrm{T} ^ {i} \in \mathrm{R} _ {\mathbf {A}} (\mathrm{G}) [ [ \mathrm{T} ] ],
\]

\[
\lambda (\mathrm{T}) = \sum_ {i = 0} ^ {\infty} \left[ \boldsymbol {\Lambda} ^ {i} (\mathrm{L}) \right] \mathrm{T} ^ {i} \in \mathrm{R} _ {\mathbf {A}} (\mathrm{G}) [ [ \mathrm{T} ] ],
\]

relations (10) are written

\begin{enumerate}
\def\labelenumi{(\arabic{enumi})}
\setcounter{enumi}{10}
\tightlist
\item
\end{enumerate}

\[
s (\mathrm{T}) \lambda (- \mathrm{T}) = 1 _ {*}.
\]

